% =====================================================================
%  3.2 Andregradsfunksjoner — teori og oppgaver
% =====================================================================

\del{1}{Graf, nullpunkter og symmetrilinje}

\begin{ordliste}
andregradsfunksjon & En funksjon $f(x) = ax^2 + bx + c$ med $a \neq 0$. & \mittsprak\\
parabel & Grafen til en andregradsfunksjon. & \mittsprak\\
symmetrilinje & Den loddrette linjen som deler parabelen i to like deler. & \mittsprak\\
bunnpunkt og toppunkt & Det laveste punktet (parabel opp) eller det høyeste punktet (parabel ned). & \mittsprak\\
ekstremalpunkt & Fellesnavn på bunnpunkt og toppunkt. & \mittsprak\\
\end{ordliste}

\begin{mustkunne}
\begin{minipage}[t]{0.57\linewidth}
\vspace{0pt}\textbf{Grafen til $f(x) = ax^2 + bx + c$.}\par
$a > 0$: parabelen vender opp (bunnpunkt). $a < 0$: den vender ned (toppunkt).\par
$f(0) = c$: parabelen skjærer $y$-aksen i $(0, c)$.\par
Nullpunktene er løsningene av $f(x) = 0$. Antallet følger av $D = b^2 - 4ac$.\medskip\par
\textbf{Symmetrilinjen.} Nullpunktene er $x_1$ og $x_2$, og $x_1 + x_2 = -\frac{b}{a}$. Symmetrilinjen ligger midt mellom dem:
\[ x = \frac{x_1 + x_2}{2} = -\frac{b}{2a} \]
Formelen inneholder ikke $c$. Å endre $c$ flytter parabelen opp eller ned, og symmetrilinjen ligger fast. Derfor gjelder formelen også uten nullpunkter.\par
\textbf{Ekstremalpunktet:} sett $x = -\frac{b}{2a}$ inn i $f$.
\end{minipage}\hfill
\begin{minipage}[t]{0.40\linewidth}
\vspace{0pt}\centering
\begin{tikzpicture}[baseline=(current bounding box.north)]
\begin{axis}[width=\linewidth, height=0.95\linewidth, axis lines=middle, xmin=-4.5, xmax=1, ymin=-2, ymax=4.5,
  xtick={-3,-1}, ytick={\empty}, tick label style={font=\scriptsize}, xlabel={$x$}, ylabel={$y$}, label style={font=\small}]
\addplot[domain=-4:0, thick, hoved, samples=40] {x^2 + 4*x + 3};
\draw[dashed, red!70!black] (axis cs:-2,-2) -- (axis cs:-2,4.3);
\addplot[only marks, mark=*, mark size=1.5pt] coordinates {(-3,0) (-1,0) (-2,-1) (0,3)};
\node[font=\scriptsize, anchor=north west, fill=white, inner sep=1pt] at (axis cs:-1.9,-1.05) {$(-2,-1)$};
\node[font=\scriptsize, anchor=south west] at (axis cs:0.05,3) {$(0,3)$};
\end{axis}
\end{tikzpicture}\par
{\footnotesize $f(x) = x^2 + 4x + 3$. Stiplet linje: $x = -2$.}
\end{minipage}
\end{mustkunne}

\begin{eksempel}[1 — graf med nullpunkter]
La $f(x) = x^2 - 2x - 3$. Finn nullpunkter, symmetrilinje, ekstremalpunkt og skjæring med $y$-aksen.

\begin{stegtabell}
$a = 1 > 0$ & Parabelen vender opp, så ekstremalpunktet er et bunnpunkt.\\
$x^2 - 2x - 3 = 0$, $x = \dfrac{2 \pm \sqrt{16}}{2}$ & Nullpunktene: abc-formelen med $D = 4 + 12 = 16$.\\
$x = -1$ eller $x = 3$ & Nullpunktene er $(-1, 0)$ og $(3, 0)$.\\
$x = \dfrac{-1 + 3}{2} = 1$ & Symmetrilinjen ligger midt mellom nullpunktene. Samme svar: $-\frac{-2}{2 \cdot 1} = 1$.\\
$f(1) = 1 - 2 - 3 = -4$ & Sett symmetrilinjens $x$ inn i $f$. Bunnpunkt $(1, -4)$.\\
$f(0) = -3$ & Skjæring med $y$-aksen: $(0, -3)$.\\
\end{stegtabell}
\end{eksempel}

\begin{eksempel}[2 — uten nullpunkter]
La $g(x) = -x^2 + 4x - 5$. Finn symmetrilinje og ekstremalpunkt.

\begin{stegtabell}
$D = 4^2 - 4\cdot(-1)\cdot(-5) = -4$ & $D < 0$: ingen nullpunkter.\\
$x = -\dfrac{4}{2\cdot(-1)} = 2$ & Symmetrilinjen finnes likevel.\\
$g(2) = -4 + 8 - 5 = -1$ & $a = -1 < 0$, så parabelen vender ned. Toppunkt $(2, -1)$.\\
Kontroll: & Toppunktet ligger under $x$-aksen, og parabelen vender ned. Da kan den ikke treffe aksen.\\
\end{stegtabell}
\end{eksempel}

\begin{vanligfeil}
\begin{itemize}[leftmargin=1.2em,itemsep=1pt,topsep=0pt]
  \item «Symmetrilinjen er $x = \frac{b}{2a}$» gjelder ikke generelt (bare når $b = 0$). Riktig er $x = -\frac{b}{2a}$. For $f(x) = x^2 + 6x$ er $x = -\frac{6}{2} = -3$, ikke $3$.
  \item «$c$ er $y$-verdien til bunnpunktet» gjelder ikke generelt (bare når $b = 0$). $c$ er skjæringen med $y$-aksen. For $x^2 - 4x + 7$ er $c = 7$, men bunnpunktet er $(2, 3)$.
  \item «Toppunkt når $a > 0$» er feil. Parabelen vender opp, så ekstremalpunktet er et bunnpunkt.
\end{itemize}
\end{vanligfeil}

\nivaahode{1}

\oppg{1.1}{1}{\uh}
Vender parabelen opp eller ned? Hvor skjærer den $y$-aksen?
\begin{tasks}(3)
  \task $f(x) = x^2 - 5x + 4$
  \task $g(x) = -2x^2 + 3x + 7$
  \task $h(x) = 3 - x^2$
\end{tasks}

\oppg{1.2}{1}{\uh}
Finn symmetrilinjen.
\begin{tasks}(3)
  \task $f(x) = x^2 - 8x + 3$
  \task $g(x) = 2x^2 + 6x - 1$
  \task $h(x) = -x^2 + 10x$
\end{tasks}

\oppg{1.3}{1}{\uh}
La $f(x) = x^2 - 4x - 5$. Fullfør.
\begin{tasks}(1)
  \task Løs $f(x) = 0$. Nullpunktene er $x = \rute$ og $x = \rute$.
  \task Symmetrilinjen er $x = \dfrac{\rute + \rute}{2} = \rute$.
  \task Bunnpunktet er $(\rute, f(\rute)) = (\rute, \rute)$.
\end{tasks}

\nivaahode{2}

\oppg{1.4}{2}{\uh}
Finn nullpunkter, symmetrilinje og ekstremalpunkt. Si om det er bunnpunkt eller toppunkt.
\begin{tasks}(1)
  \task $f(x) = x^2 + 2x - 15$
  \task $g(x) = -x^2 + 6x - 8$
  \task $h(x) = 2x^2 - 4x + 5$
\end{tasks}

\oppg{1.5}{2}{\uh}
Elin skal finne ekstremalpunktet til $f(x) = 3x^2 + 12x + 5$. Hun har gjort to feil. Finn begge, og løs riktig.
\[ x = -\frac{12}{2\cdot 3} = -2 \;\longrightarrow\; f(-2) = 3\cdot(-2)^2 + 12\cdot(-2) + 5 \;\longrightarrow\; -12 - 24 + 5 \;\longrightarrow\; -31 \]
\[ \text{Toppunkt } (-2, -31) \]

\oppg{1.6}{2}{\mh}
Tegn grafene til $f(x) = x^2 - 4x + c$ i GeoGebra for $c = -1$, $c = 4$ og $c = 7$. Bruk gjerne en glider for $c$.
\begin{tasks}(1)
  \task Hva er likt for de tre parablene, og hva er ulikt? Hvor mange nullpunkter har hver av dem?
  \task Hvor ligger symmetrilinjen? Finn bunnpunktet uttrykt ved $c$.
  \task For hvilken $c$ har parabelen ett nullpunkt?
\end{tasks}

\oppg{1.7}{2}{\uh}
Figuren viser grafen til en andregradsfunksjon $f$.
\begin{center}
\begin{tikzpicture}
\begin{axis}[width=7cm, height=5cm, axis lines=middle, xmin=-1, xmax=7, ymin=-6, ymax=5.5,
  xtick={1,...,6}, ytick={-5,...,4}, tick label style={font=\tiny}, xlabel={$x$}, ylabel={$y$}, label style={font=\small}]
\addplot[domain=-0.3:6.3, thick, hoved, samples=50] {-(x-1)*(x-5)};
\addplot[only marks, mark=*, mark size=1.5pt] coordinates {(1,0) (5,0) (3,4) (0,-5)};
\end{axis}
\end{tikzpicture}
\end{center}
\begin{tasks}(2)
  \task Les av nullpunktene.
  \task Les av symmetrilinjen.
  \task Les av toppunktet og skjæringen med $y$-aksen.
  \task Finn funksjonsuttrykket.
\end{tasks}

\nivaahode{3}

\oppg{1.8}{3}{\uh}
Er påstanden alltid, noen ganger eller aldri sann? Begrunn svaret.
\begin{tasks}(1)
  \task Parabelen til $f(x) = ax^2 + bx + c$ med $a > 0$ har to nullpunkter.
  \task Har $a$ og $c$ motsatt fortegn, har parabelen to nullpunkter.
  \task Har $a > 0$ og $c > 0$, har parabelen ingen nullpunkter.
\end{tasks}

\oppg{1.9}{3}{\uh}
En parabel $f(x) = x^2 + bx + c$ har symmetrilinje $x = -1$ og nullpunkt $x = 2$.
\begin{tasks}(1)
  \task Finn det andre nullpunktet uten å regne ut $b$ og $c$. Forklar.
  \task Finn $b$ og $c$.
\end{tasks}

\begin{sjekk}
\begin{tasks}(2)
  \task Finn symmetrilinjen til $f(x) = x^2 - 10x + 7$.
  \task Vender $f(x) = 4 - x^2$ opp eller ned? Hvor skjærer den $y$-aksen?
  \task! Forklar hvorfor symmetrilinjen ligger midt mellom nullpunktene.
\end{tasks}
\end{sjekk}

% ---------------------------------------------------------------------

\del{2}{Fullstendige kvadraters metode og toppunktform}

\begin{ordliste}
fullstendig kvadrat & Et uttrykk på formen $(x + s)^2$. Gange ut gir $x^2 + 2sx + s^2$. & \mittsprak\\
toppunktform & Funksjonen skrevet som $f(x) = a(x - p)^2 + q$. Ekstremalpunktet er $(p, q)$. & \mittsprak\\
verdimengde & Alle verdiene $f(x)$ kan få. Skrives $V_f$. & \mittsprak\\
maksimum og minimum & Den største og den minste verdien funksjonen har. & \mittsprak\\
\end{ordliste}

\begin{mustkunne}
\begin{minipage}[t]{0.50\linewidth}
\textbf{Hvorfor toppunktformen virker.}\par
Et kvadrat er aldri negativt: $(x - p)^2 \ge 0$, og lik 0 bare for $x = p$.\par
Når $a > 0$, er $f(x) = a(x - p)^2 + q \ge q$. Vi har $f(x) = q$ bare når $x = p$. Altså er $(p, q)$ bunnpunktet.\par
Når $a < 0$, snur ulikheten: $f(x) \le q$, og $(p, q)$ er toppunktet.\medskip\par
\textbf{Verdimengde:}\par
$a > 0$: $V_f = [q, \rightarrow\rangle$\par
$a < 0$: $V_f = \langle\leftarrow, q]$
\end{minipage}\hfill
\begin{minipage}[t]{0.47\linewidth}
\textbf{Fullstendige kvadraters metode.}\par
Fra 1.2: $(x + s)^2 = x^2 + 2sx + s^2$. Halvparten av tallet foran $x$ er $s$:
\[ x^2 + bx = \left(x + \tfrac{b}{2}\right)^2 - \left(\tfrac{b}{2}\right)^2 \]
Vi legger til $s^2$ og trekker det fra igjen. Da endres ikke verdien.\par
Er $a \neq 1$: sett $a$ utenfor først.
\end{minipage}
\end{mustkunne}

\begin{eksempel}[1 — når $a = 1$]
Skriv $f(x) = x^2 - 6x + 5$ på toppunktform. Finn bunnpunkt og verdimengde.

\begin{stegtabell}
$f(x) = x^2 - 6x + 9 - 9 + 5$ & Halvparten av $-6$ er $-3$, og $(-3)^2 = 9$. Legg til og trekk fra 9.\\
$f(x) = (x - 3)^2 - 4$ & $x^2 - 6x + 9 = (x - 3)^2$. Og $-9 + 5 = -4$.\\
Kontroll: & $(x-3)^2 - 4 = x^2 - 6x + 9 - 4 = x^2 - 6x + 5$. Det stemmer.\\
Bunnpunkt $(3, -4)$ & $a = 1 > 0$. Avlest som $(p, q)$.\\
$V_f = [-4, \rightarrow\rangle$ & Funksjonen er aldri mindre enn $-4$, og verdien $-4$ er med.\\
\end{stegtabell}
\end{eksempel}

\begin{eksempel}[2 — når $a \neq 1$]
Skriv $f(x) = 2x^2 + 8x + 3$ på toppunktform.

\begin{stegtabell}
$f(x) = 2(x^2 + 4x) + 3$ & Sett 2 utenfor de to første leddene.\\
$f(x) = 2\big((x + 2)^2 - 4\big) + 3$ & Inni parentesen: halvparten av 4 er 2, og $2^2 = 4$.\\
$f(x) = 2(x + 2)^2 - 8 + 3$ & Gang inn faktoren 2 med $-4$.\\
$f(x) = 2(x + 2)^2 - 5$ & $-8 + 3 = -5$.\\
Kontroll: & $2(x^2 + 4x + 4) - 5 = 2x^2 + 8x + 8 - 5 = 2x^2 + 8x + 3$. Det stemmer.\\
Bunnpunkt $(-2, -5)$, $V_f = [-5, \rightarrow\rangle$ & $a = 2 > 0$.\\
\end{stegtabell}
\end{eksempel}

\begin{vanligfeil}
\begin{itemize}[leftmargin=1.2em,itemsep=1pt,topsep=0pt]
  \item $f(x) = 2(x + 3)^2 - 1$ har bunnpunkt $(3, -1)$ er feil. Faktoren er $x - (-3)$, så bunnpunktet er $(-3, -1)$.
  \item For $2(x + 3)^2 - 1$ er $V_f = [-1, \rightarrow\rangle$, ikke $\langle -1, \rightarrow\rangle$. Verdien $-1$ er med, fordi $f(-3) = -1$.
  \item $2x^2 + 8x + 3 = 2(x + 2)^2 - 4 + 3$ er feil. Faktoren 2 må også ganges med $-4$: det gir $2(x + 2)^2 - 8 + 3$.
\end{itemize}
\end{vanligfeil}

\nivaahode{1}

\oppg{2.1}{1}{\uh}
Skriv på formen $(x - p)^2 + q$. Tallet $p$ kan være negativt.
\begin{tasks}(3)
  \task $x^2 + 2x + 5$
  \task $x^2 - 10x + 30$
  \task $x^2 - 8x + 9$
\end{tasks}

\oppg{2.2}{1}{\uh}
Les av ekstremalpunktet. Si om det er bunnpunkt eller toppunkt, og skriv verdimengden.
\begin{tasks}(3)
  \task $f(x) = (x - 2)^2 + 7$
  \task $g(x) = -(x + 1)^2 + 4$
  \task $h(x) = 3(x + 4)^2 - 2$
\end{tasks}

\oppg{2.3}{1}{\uh}
Fullfør utregningen av $f(x) = x^2 + 16x + 60$.
\[ f(x) = x^2 + 16x + \rute - \rute + 60 = (x + \rute)^2 - \rute \]
Bunnpunktet er $(\rute, \rute)$.

\nivaahode{2}

\oppg{2.4}{2}{\uh}
Skriv på toppunktform. Finn ekstremalpunktet og verdimengden.
\begin{tasks}(3)
  \task $2x^2 + 12x + 7$
  \task $3x^2 - 6x + 5$
  \task $-x^2 + 4x + 1$
\end{tasks}

\oppg{2.5}{2}{\uh}
Finn \textbf{den første feilen}. Skriv $f(x) = 2x^2 - 12x + 5$ riktig på toppunktform.
\[ 2x^2 - 12x + 5 \;\longrightarrow\; 2(x^2 - 6x) + 5 \;\longrightarrow\; 2\big((x - 3)^2 - 9\big) + 5 \;\longrightarrow\; 2(x - 3)^2 - 9 + 5 \;\longrightarrow\; 2(x - 3)^2 - 4 \]

\oppg{2.6}{2}{\uh}
En ball kastes oppover. Høyden over bakken er $h(t) = -5t^2 + 10t + 2$ meter etter $t$ sekunder.
\begin{tasks}(1)
  \task Skriv $h(t)$ på toppunktform.
  \task Hvor høyt kommer ballen, og når skjer det?
  \task Fra hvilken høyde kastes ballen?
\end{tasks}

\nivaahode{3}

\oppg{2.7}{3}{\uh}
Parabelen til en andregradsfunksjon har bunnpunkt $(2, -3)$ og går gjennom $(0, 5)$.
\begin{tasks}(1)
  \task Finn funksjonsuttrykket.
  \task Hvor mange nullpunkter har parabelen? Forklar uten abc-formelen.
\end{tasks}

\oppg{2.8}{3}{\uh}
Du har 40 m gjerde og skal lage en rektangulær hage inntil en vegg. Gjerde trengs på tre sider. Bredden vinkelrett på veggen er $x$ meter.
\begin{center}
\begin{tikzpicture}[x=1cm,y=1cm]
\draw[very thick, gray] (-0.5,2) -- (4,2);
\node[font=\small, gray, anchor=south] at (1.75,2.05) {vegg};
\draw[thick, hoved] (0,2) -- (0,0) -- (3.5,0) -- (3.5,2);
\node[font=\small, anchor=east] at (-0.1,1) {$x$};
\node[font=\small, anchor=west] at (3.6,1) {$x$};
\node[font=\small, anchor=north] at (1.75,-0.1) {$40 - 2x$};
\end{tikzpicture}
\end{center}
\begin{tasks}(1)
  \task Vis at arealet er $A(x) = -2x^2 + 40x$.
  \task Finn største areal og målene på hagen.
  \task Forklar hvorfor $x = 25$ ikke er en mulig bredde, selv om $A(25)$ kan regnes ut.
\end{tasks}

\oppg{2.9}{3}{\uh}
Er påstanden alltid, noen ganger eller aldri sann? Begrunn svaret.
\begin{tasks}(1)
  \task Når $a < 0$ og $q < 0$, skjærer grafen til $f(x) = a(x - p)^2 + q$ $x$-aksen.
  \task Grafen til $f(x) = a(x - p)^2 + q$ skjærer $x$-aksen.
  \task To andregradsfunksjoner med samme bunnpunkt er like.
\end{tasks}

\begin{sjekk}
\begin{tasks}(2)
  \task Skriv $x^2 - 12x + 40$ på formen $(x - p)^2 + q$.
  \task Finn bunnpunkt og verdimengde til $g(x) = 2(x + 1)^2 - 6$.
  \task! Forklar hvorfor $f(x) \ge q$ for $f(x) = a(x - p)^2 + q$ når $a > 0$.
\end{tasks}
\end{sjekk}
